\section{CEO 作为桥梁：研究、商业化和组织预期}

上市前后，CEO 的角色会发生变化。张鹏说，自己现在更多精力放在公司日常运营，尤其是前台、市场化和商业化。他还提到，CEO 很多时候是桥梁，是搭台子的人：让研究部门、技术部门、商业化团队能沟通，让大家发挥想象力、能力和执行力。

\subsection{研究和商业化不能两张皮}

模型公司很容易出现两张皮：研究团队关心模型能力，商业团队关心客户交付。张鹏强调，内部要不断对齐市场变化、技术研发和商业化需求。否则，研究可能做出漂亮但难落地的东西，商业也可能为了短期收入牺牲长期技术方向。

\begin{importantbox}{模型公司 CEO 的双重任务}
一边要保护长期研究和模型能力，一边要确保客户价值、收入、合规和组织效率。CEO 的核心不是替所有人做决定，而是让这些系统能互相听见。
\end{importantbox}

\subsection{几个动容瞬间}

CEO 的桥梁角色有时很抽象，所以本节用几个具体瞬间把它落地。它们说明模型公司不是只在论文和发布会里证明自己。

访谈最后，张鹏回忆几个瞬间。一个是在深圳为大客户待了小半年，最后不是空手回北京，而是带着几千万合同回去；这说明技术能被客户买单。一个是发布会上用智能体给大家发红包，虽然金额填错一位，但红包发出去了，被评论为 AI 给人类发的第一个红包。还有 GLM-4.5 开源后海外开发者和产品开始使用，甚至出现各种二次使用、蒸馏和包装。

这些瞬间共同说明一件事：模型公司的成就感不只来自 benchmark，也来自客户合同、真实使用、开发者认可和技术回响。

\begin{knowledgebox}{从 benchmark 到真实回响}
Benchmark 能证明能力，客户合同证明价值，开源使用证明生态，发布现场证明体验。模型公司要同时管理这几种反馈。
\end{knowledgebox}

\subsection{本章小结}

张鹏的 CEO 叙事并不“明星化”，更像一个桥梁角色：把研究、工程、商业、客户、治理和团队预期连接起来。这也是智谱“水泥感”的另一面。

